% ============================================================
%  Lab Project Proposal — Topology Optimisation of Planar
%  Elastic Energy-Storing Structures in a Fixed Area
%
%  Author : Heah Hung Xun (Jack)
%  Course : <fill in course code>
%  Supervisor : <fill in>
%
%  Compile with:  pdflatex  (or xelatex / lualatex if you prefer)
%  Encoding : UTF-8
% ============================================================

\documentclass[11pt,a4paper]{article}

% ---------- Page geometry ----------
\usepackage[a4paper, margin=2.5cm, headheight=14pt]{geometry}

% ---------- Fonts (your XCharter preamble) ----------
\usepackage[T1]{fontenc}
\usepackage[utf8]{inputenc}
\usepackage{XCharter}                    % body serif
\usepackage[scaled=0.96]{helvet}         % sans for headings
\usepackage[varqu,varl]{inconsolata}     % monospace
\usepackage[bigdelims,vvarbb]{newtxmath} % math to pair with XCharter
\usepackage[cal=boondoxo]{mathalpha}

% ---------- Core packages ----------
\usepackage{amsmath,amssymb,amsthm,mathtools}
\usepackage{physics}
\usepackage{siunitx}
\usepackage{booktabs}
\usepackage{array}
\usepackage{tabularx}
\usepackage{graphicx}
\usepackage{float}
\usepackage{caption}
\usepackage{subcaption}
\usepackage{enumitem}
\usepackage{microtype}
\usepackage{csquotes}

% ---------- Headings (minimal, sans, small caps) ----------
\usepackage{titlesec}
\titleformat{\section}
  {\normalfont\sffamily\large\bfseries}{\thesection}{0.7em}{}
\titleformat{\subsection}
  {\normalfont\sffamily\normalsize\bfseries}{\thesubsection}{0.6em}{}
\titleformat{\subsubsection}
  {\normalfont\sffamily\small\itshape}{\thesubsubsection}{0.5em}{}
\titlespacing{\section}{0pt}{1.4em}{0.5em}
\titlespacing{\subsection}{0pt}{1.1em}{0.3em}

% ---------- Header / footer ----------
\usepackage{fancyhdr}
\pagestyle{fancy}
\fancyhf{}
\fancyhead[L]{\footnotesize\sffamily Lab Project Proposal}
\fancyhead[R]{\footnotesize\sffamily \thepage}
\renewcommand{\headrulewidth}{0pt}

% ---------- Colors (subtle) ----------
\usepackage[dvipsnames]{xcolor}
\definecolor{accent}{HTML}{2E4057}

% ---------- Hyperlinks ----------
\usepackage[colorlinks=true,
            linkcolor=accent,
            citecolor=accent,
            urlcolor=accent]{hyperref}

% ---------- Bibliography ----------
\usepackage[backend=biber, style=numeric-comp, sorting=none]{biblatex}
\addbibresource{references.bib}

% ---------- Theorems / claim environments ----------
\newtheoremstyle{claim}
  {0.8em}{0.8em}        % space above, below
  {\itshape}             % body font
  {}                     % indent
  {\bfseries\sffamily}   % head font
  {.}{0.5em}             % punctuation after head, space
  {}                     % head spec
\theoremstyle{claim}
\newtheorem{hypothesis}{Hypothesis}

% ---------- Custom macros ----------
\newcommand{\Vf}{V_{\!f}}
\newcommand{\Um}{U/m}
\newcommand{\sy}{\sigma_{\!y}}
\newcommand{\xy}{x_{\!y}}
\newcommand{\Ur}{U_r}
\newcommand{\eff}{\text{eff}}

% ---------- Title block ----------
\title{
  \sffamily\bfseries\large
  Topology Optimisation of Planar Elastic \\
  Energy-Storing Structures in a Fixed Area
  \vspace{0.3em}\\
  \normalfont\mdseries\normalsize Lab Project Proposal
}
\author{
  \normalsize
  <Your name>\\
  \footnotesize <Student ID> \quad <Programme>\\
  \footnotesize Supervisor: <supervisor name>
}
\date{\footnotesize\today}

% ============================================================
\begin{document}
% ============================================================

\maketitle
\thispagestyle{empty}

\begin{abstract}
  \noindent
  % One short paragraph (approx. 150-200 words):
  % 1-2 sentences of motivation,
  % 1 sentence stating the two hypotheses at a high level,
  % 1-2 sentences on the rig and material,
  % 1 sentence on expected outcomes.
\end{abstract}

\vspace{1em}

% ------------------------------------------------------------
\section{Introduction and motivation}
% ------------------------------------------------------------
% Brief context: compliant mechanisms, topology optimisation
% in planar elasticity, why energy storage rather than
% compliance minimisation, and why this is an undergraduate-
% level experimental project rather than pure simulation.

% ------------------------------------------------------------
\section{Background}
% ------------------------------------------------------------

\subsection{Topology optimisation in the fixed-area setting}
% What TO is, what the design region means, SIMP in 88 lines,
% the Michell-like optimum that linear-elastic compliance
% minimisation tends towards.

\subsection{Why linear compliance minimisation is not enough}
% Cite Alacoque-James 2024, Chen-Zhang-Zhu 2018. The
% interesting results come from nonlinear / buckling-aware
% formulations; this motivates H2.

\subsection{Volume fraction as the experimental control}
% Define $V_f$, explain why fixing it decouples topology
% from material budget, and give the target mass for the
% chosen region.

% ------------------------------------------------------------
\section{Hypotheses}
% ------------------------------------------------------------

\begin{hypothesis}[Geometry screening]
  % One-paragraph claim about how structure class controls
  % specific energy under a fixed mass budget, comparing
  % cantilever beam, Archimedean spiral, honeycomb lattice,
  % and bistable arch array in the same planar region.
\end{hypothesis}

\begin{hypothesis}[Optimisation objective]
  % One-paragraph claim about buckling-induced TO versus
  % SIMP compliance TO versus the best H1 baseline. The
  % point is that the choice of TO objective dominates the
  % result, and compliance-minimising TO may sit inside the
  % Pareto frontier defined by geometric baselines.
\end{hypothesis}

% ------------------------------------------------------------
\section{Methods}
% ------------------------------------------------------------

\subsection{Design region and mass budget}
% Specify 100 x 100 x 3 mm region, $V_f = 0.30$,
% target specimen mass, and the rationale.

\subsection{Material selection}
% Short justification for PETG vs PLA, Nylon, TPU.
% Include the resilience-modulus argument U_r = sy^2/(2E)
% and the printability / hysteresis trade-off.

\subsection{Specimen fabrication}
% Printer, nozzle, layer height, wall count, infill,
% raster angle, temperature, speed, drying protocol.
% State the acceptance criterion (mass within +/- 2 percent).

\subsection{Experimental rig: PASCO Smart Cart with pulley loading}
% Describe the cart-on-track quasi-static loading rig,
% how synchronisation of $F$ and $x$ is handled, and
% why a pulley + water-cup ramp gives quasi-static
% strain rates appropriate for PETG viscoelasticity.

\begin{figure}[H]
  \centering
  % \includegraphics[width=0.85\linewidth]{figures/rig_diagram.pdf}
  \caption{Schematic of the PASCO cart + pulley loading rig.
           Placeholder -- replace with a CAD or hand-drawn figure.}
  \label{fig:rig}
\end{figure}

\subsection{Variables and controls}

\begin{table}[H]
  \centering\small
  \caption{Variables for Hypothesis 1 (geometry screening).}
  \label{tab:h1vars}
  \begin{tabularx}{\linewidth}{@{}l l X@{}}
    \toprule
    Variable & Type & Role / measurement \\
    \midrule
    $L$, $t$ & Controlled & Region dimensions \\
    $\Vf$    & Controlled & Material volume fraction \\
    Geometry class & Independent & Beam / spiral / honeycomb / bistable \\
    $x(t)$   & Controlled & PASCO encoder \\
    $F(x)$   & Dependent  & PASCO force sensor \\
    $U$      & Dependent (derived) & $\int_0^{x^\ast} F\,dx$ \\
    $\Um$    & Primary outcome & Specific stored energy \\
    $\xy$    & Dependent  & First 0.2 percent offset yield \\
    \bottomrule
  \end{tabularx}
\end{table}

\begin{table}[H]
  \centering\small
  \caption{Variables for Hypothesis 2 (optimisation objective).}
  \label{tab:h2vars}
  \begin{tabularx}{\linewidth}{@{}l l X@{}}
    \toprule
    Variable & Type & Role / measurement \\
    \midrule
    TO objective $\lambda$ & Independent & 0 = SIMP compliance, 1 = buckling-induced \\
    Filter radius $r_{\min}$ & Controlled & Manufacturability constraint \\
    $\Um$   & Dependent & Primary comparison metric \\
    $\eta_N$ & Dependent & Fractional hysteresis per cycle \\
    $k_N$   & Dependent & Loading-branch stiffness per cycle \\
    $\Delta k / k_0$ & Dependent & Damage accumulation indicator \\
    \bottomrule
  \end{tabularx}
\end{table}

% ------------------------------------------------------------
\section{Measurement protocol}
% ------------------------------------------------------------

\subsection{Phase 0: Material characterisation}
% Tensile coupons to extract $E$, $\sigma_y$, $\rho$ for
% the specific print conditions used.

\subsection{Phase 1: Specimen preparation}
% CAD verification of target volume, print parameters,
% acceptance criteria, storage conditions.

\subsection{Phase 2: Rig calibration}
% Thermal equilibration, levelling, tare, three-mass
% verification, friction baseline capture.

\subsection{Phase 3: H1 quasi-static ramp protocol}
% Load-unload cycles, recovery pauses, data logging,
% post-processing (friction subtraction, $U$ and $\xy$
% extraction).

\subsection{Phase 4: H2 cyclic loading protocol}
% Preconditioning, 100-cycle run at fixed amplitude,
% automation approach, post-processing (loop area,
% stiffness drift, plateau hysteresis).

\subsection{Phase 5: Aggregation and statistics}
% Summary table layout, ANOVA for H1, $t$-tests for H2,
% Pareto scatter construction.

% ------------------------------------------------------------
\section{Expected results}
% ------------------------------------------------------------

\subsection{Analytic predictions}
% State the closed-form predictions: Euler-Bernoulli
% cantilever, Gibson-Ashby honeycomb, Qiu-Lang-Slocum
% bistable arch, Chen-Zhang-Zhu buckling-induced TO.
% One short paragraph per model with the governing
% expression.

\subsection{Simulation predictions}
% Reference the simulation figures below. State the
% predicted specific-energy ranking and the predicted
% Pareto structure.

\begin{figure}[H]
  \centering
  \begin{subfigure}[t]{0.48\linewidth}
    % \includegraphics[width=\linewidth]{figures/fig_H1_Fx.pdf}
    \caption{H1: $F(x)$ for geometric baselines.}
    \label{fig:h1fx}
  \end{subfigure}
  \hfill
  \begin{subfigure}[t]{0.48\linewidth}
    % \includegraphics[width=\linewidth]{figures/fig_H2_Fx.pdf}
    \caption{H2: TO variants vs. best baseline.}
    \label{fig:h2fx}
  \end{subfigure}
  \caption{Expected force-displacement curves.}
  \label{fig:fx}
\end{figure}

\begin{figure}[H]
  \centering
  \begin{subfigure}[t]{0.48\linewidth}
    % \includegraphics[width=\linewidth]{figures/fig_specific.pdf}
    \caption{Specific energy ranking.}
    \label{fig:specific}
  \end{subfigure}
  \hfill
  \begin{subfigure}[t]{0.48\linewidth}
    % \includegraphics[width=\linewidth]{figures/fig_pareto.pdf}
    \caption{Efficiency-stability Pareto scatter.}
    \label{fig:pareto}
  \end{subfigure}
  \caption{Expected performance metrics.}
  \label{fig:results}
\end{figure}

\subsection{Statistical predictions}
% One-way ANOVA expected $F$-statistic for H1;
% two-sample $t$-test predictions for H2.

% ------------------------------------------------------------
\section{Risks and limitations}
% ------------------------------------------------------------
% Known confounds: print anisotropy, PETG viscoelasticity,
% cord friction, load misalignment. How each is controlled
% or measured.

% ------------------------------------------------------------
\section{Timeline and deliverables}
% ------------------------------------------------------------
% Week-by-week schedule. Deliverables: dataset,
% figures, final report.

\begin{table}[H]
  \centering\small
  \caption{Project timeline.}
  \label{tab:timeline}
  \begin{tabularx}{\linewidth}{@{}l l X@{}}
    \toprule
    Week & Phase & Tasks \\
    \midrule
    1-2  & Planning      & \\
    3-4  & Fabrication   & \\
    5-6  & H1 testing    & \\
    7-9  & H2 testing    & \\
    10   & Analysis      & \\
    11-12& Writing       & \\
    \bottomrule
  \end{tabularx}
\end{table}

% ------------------------------------------------------------
\section*{References}
% ------------------------------------------------------------
\printbibliography[heading=none]

% ============================================================
\appendix
% ============================================================

\section{PETG mechanical properties (characterisation data)}
% Table of measured $E$, $\sigma_y$, $\rho$ from coupons.

\section{CAD geometry specifications}
% Dimensioned drawings of each test geometry.

\section{Raw data and processing scripts}
% Location of CSV files, Python processing scripts,
% Capstone project files.

\end{document}
